% solution-template.tex
\documentclass[11pt,a4paper]{article}
\usepackage[utf8]{inputenc}
\usepackage[T1]{fontenc}
\usepackage[ngerman, english]{babel}
\usepackage{geometry}
\usepackage{fancyhdr}
\usepackage{hyperref}
\usepackage{amsmath}

\geometry{margin=25mm}
\setlength{\headheight}{14pt}
\pagestyle{fancy}
\fancyhf{}
\renewcommand{\headrulewidth}{0.4pt}
\lhead{\textbf{Name:} \studentname}
\rhead{\textbf{Matr.-Nr.:} \studentnumber}
\cfoot{\thepage}

% User info - bitte anpassen
\newcommand{\studentname}{Matthias Janßen}
\newcommand{\studentnumber}{1871808}
\newcommand{\course}{Grundlagen der Computergrafik}
\newcommand{\institute}{Universität Trier}
\newcommand{\sheetnumber}{02}
\newcommand{\tutor}{Jan Niclas Ruppenthal\\Prof. Dr. Stephan Diehl}

\title{\course}
\author{\studentname\\Matrikelnummer: \studentnumber\\\institute}
\date{\today}

\begin{document}

\maketitle

\begin{center}
    \textbf{Abgabe:} Übungsblatt \sheetnumber \\
    \vspace{4mm}
    \textbf{Tutor:} \tutor
\end{center}

\bigskip

\section*{Aufgaben}
% Hier beginnen deine Lösungen.

\section{Aufgabe 1}
Für den zweiten Oktanten gilt
\[
\Delta x = x_n-x_0>0,\qquad \Delta y = y_n-y_0>0,\qquad \Delta y>\Delta x
\]
und damit $a=\Delta y/\Delta x>1$.

Da die Steigung größer als 1 ist, wird in jedem Iterationsschritt $y$ erhöht.
Zwischen zwei Schritten wird entschieden, ob $x$ unverändert bleibt oder um 1
erhöht wird. Mit einer Entscheidungsvariablen $d_i$ ergibt sich die Regel
\[
d_i<0 \Rightarrow x_{i+1}=x_i,
\qquad
d_i\ge 0 \Rightarrow x_{i+1}=x_i+1.
\]

Für die rekursive Aktualisierung folgt damit
\[
d_{i+1}=
\begin{cases}
d_i+2\Delta x, & d_i<0,\\
d_i+2(\Delta x-\Delta y), & d_i\ge 0.
\end{cases}
\]

Durch Skalierung auf ganzzahlige Arithmetik erhält man als Startwert
\[
d'_0=2\Delta x-\Delta y.
\]

\textbf{Algorithmus (Bresenham, 2. Oktant):}
\begin{verbatim}
bresenham_2_oktant(x0, y0, xn, yn)
{
    dx = xn - x0;
    dy = yn - y0;

    d  = 2*dx - dy;      // Startwert d'
    x  = x0;

    for (y = y0; y <= yn; y++) {
        drawPoint(x, y);

        if (d < 0) {
            d = d + 2*dx;
        } else {
            d = d + 2*(dx - dy);
            x = x + 1;
        }
    }
}
\end{verbatim}

\section{Aufgabe 2}
% ...

\end{document}
